\documentclass{ijisa}
\usepackage[utf8]{inputenc}
\usepackage[english]{babel}



%% Title
\title{Paper Title: Preparations of Papers for the Journals of the MECS Publisher}

%% author

\author{Firstname A. Lastname}
\affil{Name of Institution/Department, City, Zip Code, Country}
\affil{Email: first.author@hostname1.org}

\author{Firstname B. Lastname }
\author{Firstname C. Lastname}
\affil{Name of Institution/Department, City, Zip Code, Country }
\affil{	  Email: {second.author, third.author}@hostname2.org}

%% Dates and header variable 
\thisvolume{XX}
\thismonth{October}
\thisyear{20XX}
\thisdoi{XX}
\thisstartpage{01}
\thisendpage{03}

\begin{document}

\abstract{This document presents the required layout of articles to be submitted for publication in print and electronic versions of the journals of MECS Publisher. The abstract should be 160 words at least. Do not use any abbreviation and equations in your abstract.}

\keywords{
	Leave one blank line after the Abstract and write your Keywords or/and Keyphrases (5-7 words or/and phrases, separated by comma).}



\maketitle
\thispagestyle{thefirstpage}
%%\thispagestyle{fancy}


\section{Introduction}

Your goal is to simulate the usual appearance of papers in a Journal of the Academy Publisher. We are requesting that you follow these guidelines as closely as possible.

\subsection{Full-Sized Camera-Ready (CR) Copy}
\textit{Paper size:} prepare your CR paper in full-size format, on A4 paper (210 x 297 mm, 8.27 x 11.69 in).

\textit{Margins:} top = 30mm (1.18 in), bottom, left and right = 20 mm (0.79 in). 

\textit{Type sizes and typefaces:} Follow the type sizes specified in Table I. As an aid in gauging type size, 1 point is about 0.35 mm. The size of the lowercase letter “j” will give the point size. Times New Roman has to be the font for main text. Paper should be single spaced.

\textit{Column width:} 82mm (3.23 in). The space between the two columns is 6mm (0.24 in). 

\textit{Paragraph indentation:} first-line 3.7 mm (0.15 in). For Abstract and Index Terms, no first-line indentation.

\textit{Alignment:} left- and right-justify your columns. Use tables and figures to adjust column length. On the last page of your paper, adjust the lengths of the columns so that they are equal. Use automatic hyphenation and check spelling. Digitize or paste down figures.

\textit{Title:} use 24-point Times New Roman font. Its paragraph description should be set so that the line spacing is single with 6-point spacing before and 6-point spacing after. Use three additional line spacings of 10 points before the beginning of the double column section, as shown above

\paragraph{Section headings:} each major section begins with a Heading in 10 point Times New Roman font centered within the column and numbered using Roman numerals (except for ACKNOWLEDGEMENT and REFERENCES), followed by a period, two spaces, and the title using an  initial capital letter for each word. The remaining letters are in SMALL CAPITALS (8 point). The paragraph description of the section heading line should be set for 12 points before and 6 points after. 

\paragraph{Subheadings:}should be 10 point, italic, left justified, and numbered with letters (A, B, …), followed by a period, two spaces, and the title using an initial capital letter for each word. The paragraph description of the subheading line should be set for 6 points before and 3 points after. 



%% table 

% Please add the following required packages to your document preamble:
% 
\begin{table}[ht]
	\centering
	\caption{Type Sizes for Camera-Ready Papers}
	\label{tab:font-size}
	\resizebox{\columnwidth}{!}{%
		\begin{tabular}{|l|l|l|l|}
			\hline
			\begin{tabular}[c]{@{}l@{}}Type\\ size\\ (pts)\end{tabular} & Regular & Bold & Italic \\ \hline
			6 & Table captions, table superscripts &  &  \\ \hline
			8 & \begin{tabular}[c]{@{}l@{}}Section tiles, tables, tables names, \\ first letter in table captions, figure \\ captions, footnotes, text subscripts\\  and superscripts\end{tabular} &  &  \\ \hline
			9 & References, authors' biographies &  &  \\ \hline
			10 & \begin{tabular}[c]{@{}l@{}}Authors' affiliations, main text, \\ equations, first letter in section\\ title\end{tabular} &  & Subheading \\ \hline
			11 &  & \begin{tabular}[c]{@{}l@{}}Author's\\ names\end{tabular} &  \\ \hline
			24 & Paper title &  &  \\ \hline
		\end{tabular}%
	}
\end{table}

\subsection{PDF Creation}
The PDF document should be sent as an open file, i.e. without any data protection. 

Please do not use the Adobe Acrobat PDFWriter to generate the PDF file. Use the Adobe Acrobat Distiller instead, which is contained in the same package as the Acrobat PDFWriter. 

Make sure that you have used Type 1 or True Type Fonts (check with the Acrobat Reader or Acrobat Writer by clicking on File>Document Properties>Fonts to see the list of fonts and their type used in the PDF document). 

As always with a conversion to PDF, authors should very carefully check a printed copy. 


\section{HELPFUL HINTS}

\subsection{Figures and Tables}

Position figures and tables at the tops and bottoms of columns. Avoid placing them in the middle of columns. Large figures and tables may span across both columns. Figure captions should be centered below the figures; table captions should be centered above. Avoid placing figures and tables before their first mention in the text. Use the abbreviation “Fig. 1,” even at the beginning of a sentence. 


\begin{figure}[htb]
\centering
\includegraphics[width=.8\columnwidth]{figure1} 
\caption{Note how the caption is centered in the column} \label{fig:fig1}
\end{figure}


\subsection{References}
Number citations consecutively in square brackets \cite{eason1955certain}. No punctuation follows the bracket \cite{maxwell1881treatise}. Use “Ref. \cite{jacobs1963fine}” or “Reference \cite{jacobs1963fine}” at the beginning of a sentence: 
Give all authors’ names; use “et al.” if there are six authors or more. Papers that have not been published, even if they have been submitted for publication, should be cited as “unpublished” [4]. Papers that have been accepted for publication should be cited as “in press” [5]. In a paper title, capitalize the first word and all other words except for conjunctions, prepositions less than seven letters, and prepositional phrases.
For papers published in translated journals, first give the English citation, then the original foreign-language citation [\cite{yorozu1987electron}.
For on-line references a URL and time accessed must be given. 
At the end of each reference, give the DOI (Digital Object Identifier) number as long as available, in the format as “doi:10.1518/hfes.2006.27224”
\subsection{Footnotes}
Number footnotes separately in superscripts 1, 2, …. Place the actual footnote at the bottom of the column in which it was cited, as in this column. See first page footnote for an example. 
Dates of manuscript submission, revision and acceptance should be included in the first page footnote. Remove the first page footnote if you don’t have any information there.
\subsection{Abbreviations and Acronyms}
Define abbreviations and acronyms the first time they are used in the text, even after they have been defined in the abstract. Do not use abbreviations in the title unless they are unavoidable.
\subsection{Equations}
Equations should be centered in the column. The paragraph description of the line containing the equation should be set for 6 points before and 6 points after. Number equations consecutively with equation numbers in parentheses flush with the right margin, as in (1). Italicize Roman symbols for quantities and variables, but not Greek symbols. Punctuate equations with commas or periods when they are part of a sentence, as in	(1)
Symbols in your equation should be defined before the equation appears or immediately following. Use “(1),” not “Eq. (1)” or “equation (1),” except at the beginning of a sentence: "Equation (1) is ..."


\subsection{Other Recommendations}
Use either SI (MKS) or CGS as primary units. (SI units are encouraged.) If your native language is not English, try to get a native English-speaking colleague to proofread your paper. Do not add page numbers.
\subsection{A Quick Checklist} 
\begin{itemize}
	\item Paper size=A4; Margins: top=3cm, bottom=left=right=2cm; Column spacing=0.6cm
	\item For the whole document (“Ctrl-A” to select the whole document), Font Type=Times New Roman, do NOT use any Asian font type like SimSun in formulas, section numbers (III, IV, V, ...), list numbers (1), 2), (1), (2), ...), or punctuation marks (“,”, “.”, “:”, “;”, “(“, “)”, ...). Check Word Count (on the status bar at the bottom of the window) to ensure the number of Asian Characters (including textboxes and footnotes) is 0
	\item 	In Paragraph settings for the whole document ( “Ctrl-A” to select the whole document), Line spacing must be "Single", "Snap to grid" must NOT be checked. 
	\item In Paragraph settings for main text except section titles, Indentation left=right=0, first line=0.37cm; Spacing before=after=0, not blank line between paragraphs
	
	\item Title and authors: font style=regular NOT bold NOT italic; font size for title is 24 with 6 spacing before \& after, for authors names font size is 11, affiliations font size is 10
	
	\item Biographies: it is strongly recommended adding for each author a short bio to the end of the paper.
	
\end{itemize}


\section{APPENDIX A APPENDIX TITLE}
Appendixes, if needed, are numbered by A, B, C... Use two spaces before APPENDIX TITLE.


\section{ACKNOWLEDGMENT}
	The authors wish to thank A, B, C. This work was supported in part by a grant from XYZ.	

%% include bib file
\bibliographystyle{unsrt}
\bibliography{ijisa_sample}

%% author biography 
\begin{wrapfigure}{l}{20mm} 
    \includegraphics[width=1in,height=1.10in,keepaspectratio]{author1.png}
  \end{wrapfigure} \par
 \textbf{Author A} and the other authors may include biographies and photographs at the end of regular papers. Photographs, if provided, should be cropped into 26mm in width and 32mm in height. The first paragraph may contain a place and/or date of birth (list place, then date). Next, the author’s educational background is listed. The degrees should be listed with type of degree in what field, which institution, city, state or country, and year degree was earned. The author’s major field of study should be lower-cased.
The second paragraph uses the pronoun of the person (he or she) and not the author’s last name. It lists military and work experience, including summer and fellowship jobs. Job titles are capitalized. The current job must have a location; previous positions may be listed without one. Information concerning previous publications may be included. Try not to list more than three books or published articles. The format for listing publishers of a book within the biography is: title of book (city, state: publisher name, year) similar to a reference. Current and previous research interests end the paragraph.
The third paragraph begins with the author’s title and last name (e.g., Dr. Smith, Prof. Jones, Mr. Kajor, Ms. Hunter). List any memberships in professional societies like the IEEE. Finally, list any awards and work for professional committees and publications. Personal hobbies should not be included in the biography. \par \par

\begin{wrapfigure}{l}{20mm} 
\vspace{-5pt}
    \includegraphics[width=1in,height=1in,clip, keepaspectratio]{author1.png}
  \end{wrapfigure}\par
 \textbf{Author B} and the other authors may include biographies and photographs at the end of regular papers. Photographs, if provided, should be cropped into 26mm in width and 32mm in height. The first paragraph may contain a place and/or date of birth (list place, then date). Next, the author’s educational background is listed. The degrees should be listed with type of degree in what field, which institution, city, state or country, and year degree was earned. The author’s major field of study should be lower-cased. \par
 
 
\begin{wrapfigure}{l}{0.15\textwidth}
  \vspace{-15pt}
  \begin{center}
    \includegraphics[width=0.15\textwidth]{author1}
  \end{center}
  \vspace{-20pt}
  \vspace{-10pt}
\end{wrapfigure}
 \textbf{Author C} includes the biography here. \par

\end{document}